% !TEX program = xelatex
\documentclass[11pt,a4paper]{article}

\usepackage{fontspec}
\usepackage{polyglossia}
\usepackage{amsmath,amssymb,amsfonts}
\usepackage{xcolor}
\usepackage{graphicx}
\usepackage{geometry}

\geometry{left=2cm,right=2cm,top=2cm,bottom=2cm}

\setdefaultlanguage{persian}
\setotherlanguage{english}
\setmainfont{XB Niloofar}

\newcommand{\en}[1]{\textenglish{#1}}
\newcommand{\separator}{\vspace{1em}\noindent\rule{\linewidth}{0.4pt}\vspace{1em}}
\newcommand{\thickseparator}{\vspace{1em}\noindent\rule{\linewidth}{1pt}\vspace{1em}}

\title{\textbf{پاسخ سوال ۳} \\[0.4cm] \small مدل Skip-Gram با نمونه‌گیری منفی (SGNS)}
\author{محمد باغبان زاده \\ شماره دانشجویی: شناسهٔ دانشجو}
\date{}
\linespread{1.35}

\begin{document}

\maketitle
\thickseparator

\section{مفاهیم پایه: مدل Skip-Gram با نمونه‌گیری منفی (SGNS)}

\subsection{معرفی مدل Word2vec و Skip-Gram}
در بخش \textbf{۵.۵ کتاب}، مدل \textbf{word2vec} به عنوان روشی برای تولید بردارهای کلمه (embeddings) با ابعاد کم و متراکم معرفی شده است. برخلاف روش‌های مبتنی بر شمارش ساده (بخش ۵.۳) که بردارهای بسیار طویل و پراکنده (sparse) تولید می‌کنند، word2vec بردارهایی با ابعاد \(d\) (معمولاً بین ۵۰ تا ۱۰۰۰) تولید می‌کند که مقادیر آنها حقیقی و می‌توانند منفی باشند.

الگوریتم \textbf{Skip-Gram} یکی از دو الگوریتم موجود در word2vec است. ایده اصلی آن به این صورت است:

\begin{itemize}
    \item به ازای هر کلمه هدف \(w\)، کلمات همسایه آن در پنجره‌ای به طول \(L\) را به عنوان \textbf{نمونه‌های مثبت} در نظر می‌گیریم.
    \item کلمات تصادفی از دایره واژگان را به عنوان \textbf{نمونه‌های منفی} انتخاب می‌کنیم.
    \item یک طبقه‌بند Logistic Regression آموزش می‌دهیم تا بین این دو دسته تمایز قائل شود.
    \item وزن‌های یادگرفته شده توسط این طبقه‌بند، همان بردارهای کلمه (embeddings) خواهند بود.
\end{itemize}

\subsection{نمونه‌گیری منفی \en{(Negative Sampling)}}
در بخش \textbf{۵.۵.۲ کتاب}، روش \textbf{نمونه‌گیری منفی \en{(Negative Sampling)}} توضیح داده شده است. به ازای هر نمونه مثبت \((w, c_{pos})\)، تعداد \(k\) نمونه منفی \((w, c_{neg})\) از دایره واژگان نمونه‌برداری می‌شود. احتمال انتخاب هر کلمه به عنوان نمونه منفی، متناسب با بسامد آن کلمه به توان \(\alpha = 0.75\) است:

\[
P_\alpha(w) = \frac{\text{count}(w)^\alpha}{\sum_{w'} \text{count}(w')^\alpha}
\]

\subsection{تابع هزینه \en{(Loss Function)}}
طبق \textbf{معادله \en{۵-۲۱} کتاب}، تابع هزینه برای یک جفت \((w, c_{pos})\) و \(k\) نمونه منفی \(c_{neg_1}, ..., c_{neg_k}\) به صورت زیر تعریف می‌شود:

\[
L(w, c_{pos}, c_{neg_1..k}) = -\log P(+|w, c_{pos}) - \sum_{i=1}^{k} \log P(-|w, c_{neg_i})
\]

با جایگذاری تعاریف \(P(+|w,c) = \sigma(c \cdot w)\) و \(P(-|w,c) = \sigma(-c \cdot w)\):

\[
\boxed{L = -\log \sigma(c_{pos} \cdot w) - \sum_{i=1}^{k} \log \sigma(-c_{neg_i} \cdot w)}
\]

هدف الگوریتم، کمینه کردن این تابع هزینه است؛ یعنی بیشینه کردن تشابه (دترمین داخلی) بین بردار کلمه هدف و بردار کلمات زمینه واقعی، و کمینه کردن تشابه بین بردار کلمه هدف و بردار نمونه‌های منفی.

\thickseparator

\section{حل گام‌به‌گام مسأله}

\subsection{داده‌های مسأله}
\begin{align*}
\text{:بردار کلمه هدف \en{(Target word)}} \quad v_w &= [0.4, \; 0.3] \\
\text{:بردار کلمه زمینه واقعی \en{(True Context)}} \quad v_c &= [0.5, \; 0.5] \\
\text{:نمونه منفی ۱} \quad v_{n_1} &= [0.2, \; 0.1] \\
\text{:نمونه منفی ۲} \quad v_{n_2} &= [0.1, \; -0.3] \\
\text{:تعداد نمونه‌های منفی} \quad k &= 2
\end{align*}

\subsection{گام ۱: محاسبه دترمین داخلی \en{(Dot Product)} برای هر جفت}
دترمین داخلی دو بردار \(a = [a_1, a_2]\) و \(b = [b_1, b_2]\) به صورت \(a \cdot b = a_1b_1 + a_2b_2\) محاسبه می‌شود.

\begin{align*}
v_c \cdot v_w &= (0.5 \times 0.4) + (0.5 \times 0.3) = 0.20 + 0.15 = 0.35 \\[4pt]
v_{n_1} \cdot v_w &= (0.2 \times 0.4) + (0.1 \times 0.3) = 0.08 + 0.03 = 0.11 \\[4pt]
v_{n_2} \cdot v_w &= (0.1 \times 0.4) + (-0.3 \times 0.3) = 0.04 - 0.09 = -0.05
\end{align*}

\subsection{گام ۲: اعمال تابع سیگموئید \(\sigma(x) = \frac{1}{1+e^{-x}}\)}
برای محاسبه \(\sigma(c_{pos} \cdot w) = \sigma(0.35)\) و \(\sigma(-c_{neg} \cdot w)\) به مقادیر زیر نیاز داریم:

ابتدا \(e^{-0.35}\) و \(e^{0.35}\) را تقریب می‌زنیم (\(e^{0.35} \approx 1.419\) و \(e^{-0.35} \approx 0.7047\)):

\[
\sigma(0.35) = \frac{1}{1 + e^{-0.35}} = \frac{1}{1 + 0.7047} = \frac{1}{1.7047} \approx 0.5866
\]

برای \(\sigma(-v_{n_1} \cdot v_w) = \sigma(-0.11)\) و \(\sigma(-v_{n_2} \cdot v_w) = \sigma(0.05)\):

\[
\sigma(-0.11) = \frac{1}{1 + e^{0.11}} \quad,\quad e^{0.11} \approx 1.1163 \quad,\quad \sigma(-0.11) = \frac{1}{2.1163} \approx 0.4725
\]

\[
\sigma(0.05) = \frac{1}{1 + e^{-0.05}} \quad,\quad e^{-0.05} \approx 0.9512 \quad,\quad \sigma(0.05) = \frac{1}{1.9512} \approx 0.5125
\]

\subsection{گام ۳: محاسبه لگاریتم طبیعی \en{(Natural Logarithm)}}
با استفاده از \(\log\) طبیعی (مبنای \(e\)):

\[
\log \sigma(0.35) = \log(0.5866) \approx -0.5335
\]

\[
\log \sigma(-0.11) = \log(0.4725) \approx -0.7495
\]

\[
\log \sigma(0.05) = \log(0.5125) \approx -0.6685
\]

\subsection{گام ۴: محاسبه نهایی \en{Loss} با استفاده از فرمول معادله \en{۵-۲۱}}
\begin{align*}
L &= -\log \sigma(c_{pos} \cdot w) - \sum_{i=1}^{k} \log \sigma(-c_{neg_i} \cdot w) \\
&= -(\log \sigma(0.35)) - \big[ \log \sigma(-0.11) + \log \sigma(0.05) \big] \\
&= -(-0.5335) - \big[ (-0.7495) + (-0.6685) \big] \\
&= 0.5335 - \big[ -1.4180 \big] \\
&= 0.5335 + 1.4180 \\
&= 1.9515
\end{align*}

\separator
\section{جمع‌بندی و پاسخ نهایی}

\begin{itemize}
    \item \textbf{دترمین داخلی‌ها:} \(v_c \cdot v_w = 0.35\)، \(v_{n_1} \cdot v_w = 0.11\)، \(v_{n_2} \cdot v_w = -0.05\)
    \item \textbf{مقادیر سیگموئید:} \(\sigma(0.35) \approx 0.5866\)، \(\sigma(-0.11) \approx 0.4725\)، \(\sigma(0.05) \approx 0.5125\)
    \item \textbf{لگاریتم مقادیر:} \(\log\sigma(0.35) \approx -0.5335\)، \(\log\sigma(-0.11) \approx -0.7495\)، \(\log\sigma(0.05) \approx -0.6685\)
\end{itemize}

\thickseparator
\begin{center}
\fbox{\parbox{0.85\linewidth}{\centering
\textbf{پاسخ نهایی:} \\
مقدار \en{Loss} برای این نمونه آموزشی برابر است با:
\[
\boxed{L \approx 1.9515}
\]
}}
\end{center}

\thickseparator

\section{نکته: ارتباط با فرمول کتاب}
طبق \textbf{معادله \en{۵-۲۱}} در کتاب:

\[
L(w, c_{pos}, c_{neg_1..k}) = -\log \sigma(c_{pos} \cdot w) - \sum_{i=1}^{k} \log \sigma(-c_{neg_i} \cdot w)
\]

همانطور که مشاهده می‌شود، جمله اول (\(-\log \sigma(c_{pos} \cdot w)\)) سعی می‌کند تشابه بین \(v_w\) و \(v_c\) را بیشینه کند (زیرا هرچه این تشابه بیشتر باشد، \(\sigma\) به ۱ نزدیک‌تر و \(-\log\sigma\) به صفر نزدیک‌تر می‌شود). جمله دوم (\(-\sum \log \sigma(-c_{neg} \cdot w)\)) سعی می‌کند تشابه بین \(v_w\) و نمونه‌های منفی را کمینه کند (زیرا هرچه این تشابه کمتر باشد، \(-c_{neg} \cdot w\) بزرگتر، \(\sigma\) به ۱ نزدیک‌تر و \(-\log\sigma\) به صفر نزدیک‌تر می‌شود). مقدار محاسبه‌شده \(L \approx 1.95\) نشان‌دهنده این است که مدل هنوز فاصله زیادی با وضعیت بهینه دارد.

\thickseparator

\end{document}